\section{Chains}\label{sec:chains}
% 2026-09-21 house shape: opener with theorem preview; definition -> announced, named theorem, proof <= 8 sentences
% -> example on Figure 1 -> remark; lemma; proposition; twins; how many chains.  Cadence pass: short sentences.

% 2026-09-22 cold-read fixes: the opener says what a record and the partition are; the definition reads once; the
% theorem is named for what it proves; nucleus variables renamed K to keep N for the example's community.
This section answers the first difficulty of Section~\ref{sec:intro}: which vertices an index can store as one group, with one record for all of them, for every size, every level and every query.
%
% 2026-09-22 compaction: two sentences.
The own nodes alone decide hierarchy equivalence (Theorem~\ref{thm:chains}), and its classes, the chains, form the coarsest partition of the vertices that loses no answer (Proposition~\ref{prop:coarsest}).

\stitle{Hierarchy equivalence.}
Two vertices can share one record when no query tells them apart.
\begin{definition}[hierarchy equivalence]\label{def:equiv}
Two vertices $u$ and $v$ are \textit{hierarchy-equivalent} if, for every size $s$, their core values agree, $\core{s}{u}=\core{s}{v}$, and for every level $1\le k\le\core{s}{u}$ the two vertices lie in the same $\Nuc{s}{k}$.
\end{definition}
% 2026-09-22 filler pass: two commentary sentences after the definition removed (the theorem states what the definition reduces to).
\begin{definition}[trajectory, chain]\label{def:chain}
For a vertex $v$ in some edge, that is, with $\omg{v}\ge2$, the \textit{trajectory} of $v$ is the tuple of its own nodes over the sizes $2$ to $\omg{v}$.
%
We write $\traj{v}=(\own{2}{v},\dots,\own{\omg{v}}{v})$.
%
The vertices with equal trajectories form a \textit{chain}, one chain per distinct trajectory, and the isolated vertices, those in no edge, form one more chain.
\end{definition}
\begin{theorem}[own nodes decide hierarchy equivalence]\label{thm:chains}
Two vertices $u$ and $v$ with $\omg{u},\omg{v}\ge2$ are hierarchy-equivalent if and only if $\traj{u}=\traj{v}$.
\end{theorem}
\begin{proof}
($\Leftarrow$) Equal trajectories have equal length, which gives $\omg{u}=\omg{v}$, and both core values are $0$ above that size (Section~\ref{sec:prelim}).
%
For $2\le s\le\omg{u}$ the two own nodes are equal, and the top of an own node is the core value of its vertex (Section~\ref{sec:strees}), so $\core{s}{u}=\core{s}{v}$.
%
For $k\le\core{s}{u}$, the $\Nuc{s}{k}$ containing $u$ is the node nearest the root with top at least $k$ on the path from $\own{s}{u}$ to its root (Section~\ref{sec:strees}).
%
Since $\own{s}{v}=\own{s}{u}$, the path and the node are the same for $v$.
%
($\Rightarrow$) Equal core values at every size give $\omg{u}=\omg{v}$, since $\omg{v}$ is the last size with a nonzero value.
%
For $2\le s\le\omg{u}$, let $k=\core{s}{u}=\core{s}{v}$.
%
The $\Nuc{s}{k}$ containing $u$ is $\own{s}{u}$ by definition.
%
It contains $v$ by hierarchy equivalence, and it is then the $\Nuc{s}{\core{s}{v}}$ containing $v$, that is $\own{s}{v}$.
\end{proof}
\begin{example}\label{ex:chains}
% 2026-09-22 compaction: four chains in one sentence (make_concept.py: chains {b}, {a}, {u,w}, {x}).
Consider Figure~\ref{fig:example} with the own nodes of Example~\ref{ex:tree}.
%
The trajectories give four chains: $\{b_1,b_2,b_3\}$ with (root, $N$, $B$) over sizes $2$ to $4$, $\{a_1,\dots,a_5\}$ with $(A,A,A,A)$ over sizes $2$ to $5$, $\{u_1,\dots,w_3\}$ with $(\{x,u_1,\dots,w_3\},\ N)$ over sizes $2$ and $3$, and $\{x\}$ with $(\{x,u_1,\dots,w_3\},\ N,\ B)$.
%
The vertex $a_5$ has one more edge than $a_1$ but the same own node at every size, and no query separates them.
%
The vertex $x$ shares the own nodes of $u_1$ at sizes $2$ and $3$ and separates from it at size $4$, where $x$ lies in the $4$-clique $B$ and $u_1$ in none.
\end{example}

\stitle{Communities and chains.}
Every community is a union of whole chains, and every partition with this property has each of its groups inside one chain.
\begin{lemma}[communities are unions of chains]\label{lem:union}
If $u$ lies in an $\Nuc{s}{k}$ $K$ with $k\ge1$ and $w$ is in the chain of $u$, then $w\in K$.
\end{lemma}
\begin{proof}
% 2026-09-22 cold read: the proof no longer derives an unused equality; the isolated-chain case is named.
Since $k\ge1$, $u$ lies in an $s$-clique, so its chain is the class of its trajectory, and $\traj{w}=\traj{u}$ gives $\own{s}{w}=\own{s}{u}$.
%
Since $u$ lies in $K$, $\core{s}{u}\ge k$.
%
The node $\own{s}{u}$ is a core at the level $\core{s}{u}\ge k$ that contains $u$, so it lies inside the $\Nuc{s}{k}$ containing $u$, namely $K$, by the nesting of Section~\ref{sec:strees}.
%
The node $\own{s}{u}=\own{s}{w}$ contains $w$, and hence $w\in K$.
\end{proof}
\begin{proposition}[the coarsest lossless partition]\label{prop:coarsest}
Call a partition of $V$ \textit{lossless} if every community of every size and level is a union of its blocks, the sets of the partition.
%
The chain partition is lossless, and every lossless partition is at least as fine: each of its blocks lies inside one chain.
\end{proposition}
\begin{proof}
The chain partition is lossless by Lemma~\ref{lem:union}.
%
% 2026-09-22 compaction: six sentences.
Let $P$ be lossless, and let $u$ and $v$ share a block of $P$.
%
We show that they share a chain.
%
If $\core{s}{u}>\core{s}{v}$ at some size, the community of $u$ at $(s,\core{s}{u})$ contains $u$ and not $v$, since every vertex of a core at that level has at least that core value.
%
It is then not a union of blocks of $P$.
%
The core values therefore agree at every size.
%
If $u$ is isolated, so is $v$, and both lie in the chain of isolated vertices.
%
Otherwise, if $\own{s}{u}\ne\own{s}{v}$ at a size where both have value $k\ge1$, these two $(s,k)$-cores are disjoint (Section~\ref{sec:prelim}), and the community of $u$ at $(s,k)$ again contains $u$ and not $v$.
%
Hence the own nodes agree at every size, $\traj{u}=\traj{v}$, and $u$ and $v$ share a chain.
\end{proof}
% 2026-09-22 filler pass: the remark after the proposition restated it; removed.

\stitle{Twins.}
% 2026-09-22: the twins lemma folded into the definition (its proof is one sentence); Exp-3 uses twin classes as its baseline.
Two vertices $u$ and $w$ are \textit{twins} when their closed neighbourhoods, the neighbours together with the vertex itself, are the same set, and a twin class is a maximal set of pairwise twins.
%
Twins lie in one chain.
%
Swapping them maps $G$ onto itself and every core containing one of them onto a core containing the other.
%
The two cores are the same, since they share a vertex or are the pair itself.
%
Twins are therefore hierarchy-equivalent, and they are adjacent, so Theorem~\ref{thm:chains} puts them in one chain.
%
Chains are therefore at most as many as twin classes, and on the graphs of Section~\ref{sec:exp} they are far fewer.
%
Twin classes have $1.03$ to $1.28$ vertices on average, chains $1.5$ to $154$ (Table~\ref{tab:size}): \dblp has $13{,}459$ chains for $317{,}080$ vertices ($24$ per chain) and \webuk $839$ chains for $129{,}632$ vertices ($154$ per chain).

